\providecommand{\rank}{\operatorname{r}}
\providecommand{\RR}{\mathbb R}
\providecommand{\CC}{\mathbb C}
\providecommand{\KK}{\mathbb K}

% ---------- Exercise 83 ----------
\setcounter{enumi}{82}
\item {\bfseries\boldmath
Determine whether each of the following systems of vectors is a basis of $\RR^3$.
}

\begin{enumerate}[label=\alph*)]

\item $\bigl((0,0,1),(1,1,0),(1,1,1)\bigr)$.

\begin{answer}
No. In the order given, $v_3=v_1+v_2$, so the system is dependent.
\end{answer}

\item $\bigl((1,2,-1),(-1,1,1),(2,-8,-2)\bigr)$.

\begin{answer}
No. For $v_1=(1,2,-1)$, $v_2=(-1,1,1)$, $v_3=(2,-8,-2)$,
we have $v_3=-2v_1-4v_2$.
\end{answer}

\item $\bigl((1,-1,1),(1,1,7),(1,1,-1)\bigr)$.

\begin{answer}
Yes. The determinant of the column matrix is
\[
 \det\begin{bmatrix}1&1&1\\-1&1&1\\1&7&-1\end{bmatrix}=-16\ne0.
\]
The three vectors are independent and generate $\RR^3$.
\end{answer}

\item $\bigl((1,0,0),(0,1,1)\bigr)$.

\begin{answer}
No. Every combination of $(1,0,0)$ and $(0,1,1)$ satisfies $y=z$;
for example $(0,0,1)$ is not generated.
\end{answer}

\item $\bigl((0,1,1),(0,1,0),(1,1,1),(2,1,3)\bigr)$.

\begin{answer}
No. With the vectors in the listed order,
$v_4=v_1-2v_2+2v_3$. Thus the four-vector system is dependent.
\end{answer}
\end{enumerate}
